\documentclass[12pt,a4paper]{article}

% Encoding and fonts
\usepackage[utf8]{inputenc}
\usepackage[T1]{fontenc}
\usepackage{lmodern}
\usepackage{textcomp}

% Page layout
\usepackage[top=1in, bottom=1in, left=1in, right=1in]{geometry}

% Typography
\usepackage{microtype}
\usepackage{parskip}

% Language
\usepackage[english]{babel}

% Headers and footers
\usepackage{fancyhdr}
\pagestyle{fancy}
\fancyhf{}
\fancyhead[L]{\leftmark}
\fancyhead[R]{\thepage}
\fancyfoot[C]{\thepage}
\renewcommand{\headrulewidth}{0.4pt}
\renewcommand{\footrulewidth}{0pt}

% Section styling
\usepackage{titlesec}
\titleformat{\section}{\Large\bfseries}{\thesection}{1em}{}
\titleformat{\subsection}{\large\bfseries}{\thesubsection}{1em}{}
\titleformat{\subsubsection}{\normalsize\bfseries}{\thesubsubsection}{1em}{}
\titlespacing*{\section}{0pt}{2.5ex plus 1ex minus .2ex}{1.5ex plus .2ex}
\titlespacing*{\subsection}{0pt}{2ex plus 1ex minus .2ex}{1ex plus .2ex}
\titlespacing*{\subsubsection}{0pt}{1.5ex plus 1ex minus .2ex}{0.8ex plus .2ex}

% List formatting
\usepackage{enumitem}
\setlist[itemize]{topsep=0.5ex, itemsep=0.25ex, parsep=0pt}
\setlist[enumerate]{topsep=0.5ex, itemsep=0.25ex, parsep=0pt}

% Essential packages
\usepackage{graphicx}
\usepackage[table]{xcolor}
\usepackage{booktabs}
\usepackage{array}
\usepackage{tabularx}
\usepackage{longtable}
\usepackage{amsmath}
\usepackage{amssymb}
\usepackage[normalem]{ulem}
\rowcolors{2}{gray!10}{white}
\renewcommand{\arraystretch}{1.3}

% Cross-references
\usepackage{hyperref}
\usepackage{cleveref}

% Custom commands
\newcommand{\pilcrow}{\S}

% Hyperref setup
\hypersetup{
    colorlinks=true,
    linkcolor=blue,
    filecolor=magenta,
    urlcolor=cyan,
}

% Code listings
\usepackage{listings}
\definecolor{codebg}{RGB}{248,248,248}
\definecolor{codeframe}{RGB}{200,200,200}
\definecolor{codekeyword}{RGB}{0,0,180}
\definecolor{codecomment}{RGB}{0,128,0}
\definecolor{codestring}{RGB}{163,21,21}
\lstset{
    basicstyle=\ttfamily\small,
    breaklines=true,
    breakatwhitespace=false,
    frame=single,
    framerule=0.5pt,
    rulecolor=\color{codeframe},
    backgroundcolor=\color{codebg},
    keepspaces=true,
    showstringspaces=false,
    tabsize=4,
    xleftmargin=1em,
    xrightmargin=1em,
    aboveskip=1em,
    belowskip=1em,
    keywordstyle=\color{codekeyword}\bfseries,
    commentstyle=\color{codecomment}\itshape,
    stringstyle=\color{codestring},
    literate={_}{{\textunderscore}}1
             {-}{{\textminus}}1
             {<}{{\textless}}1
             {>}{{\textgreater}}1
             {~}{{\textasciitilde}}1
             {^}{{\textasciicircum}}1
}

% YAML language definition for listings
\lstdefinelanguage{YAML}{
    keywords={true,false,null,y,n},
    keywordstyle=\color{blue}\bfseries,
    sensitive=false,
    comment=[l]{\#},
    morecomment=[s]{/*}{*/},
    commentstyle=\color{gray}\ttfamily,
    stringstyle=\color{red}\ttfamily,
    morestring=[b]',
    morestring=[b]',
}

% TOML language definition for listings
\lstdefinelanguage{TOML}{
    keywords={true,false},
    keywordstyle=\color{blue}\bfseries,
    sensitive=true,
    comment=[l]{\#},
    commentstyle=\color{gray}\ttfamily,
    stringstyle=\color{red}\ttfamily,
    morestring=[b]',
    morestring=[b]',
}

% Dockerfile language definition for listings
\lstdefinelanguage{Dockerfile}{
    keywords={FROM,RUN,CMD,LABEL,EXPOSE,ENV,ADD,COPY,ENTRYPOINT,VOLUME,USER,WORKDIR,ARG,ONBUILD,STOPSIGNAL,HEALTHCHECK,SHELL},
    keywordstyle=\color{blue}\bfseries,
    sensitive=false,
    comment=[l]{\#},
    commentstyle=\color{gray}\ttfamily,
    stringstyle=\color{red}\ttfamily,
    morestring=[b]',
    morestring=[b]',
}

% JSON language definition for listings
\lstdefinelanguage{JSON}{
    keywords={true,false,null},
    keywordstyle=\color{blue}\bfseries,
    sensitive=true,
    commentstyle=\color{gray}\ttfamily,
    stringstyle=\color{red}\ttfamily,
    morestring=[b]',
}

% Code listing float environment
\usepackage{float}
\newfloat{listing}{htbp}{lol}[section]
\floatname{listing}{Listing}

% Admonition boxes
\usepackage{tcolorbox}
\tcbuselibrary{skins,breakable}

% Admonition style definitions
\newtcolorbox{admonitionbox}[2][]{
    enhanced,
    breakable,
    colback=#2!5,
    colframe=#2!80!black,
    fonttitle=\bfseries,
    title=#1,
    left=2mm,
    right=2mm,
}

% Synthon tuple boxes
\newtcolorbox{synthonbox}{
    enhanced,
    colback=blue!5,
    colframe=blue!75!black,
    fontupper=\large,
    center upper,
    boxrule=1pt,
    arc=4pt,
}

\begin{document}

\title{Analytical Foundations of the MillenniumAnkh Framework and ZFCₜ}
\date{\today}

\maketitle

\subsection{I. Executive Summary}

The MillenniumAnkh library targets the formalization of the remaining seven Millennium Prize Problems in Lean 4. Rather than claiming proofs of open conjectures, the framework establishes a \textbf{Machine-Checked Barrier Taxonomy} that classifies the epistemic and ontological obstacles preventing their resolution. Central to this taxonomy is the distinction between \texttt{MathlibGap} (formalization deficit), \texttt{OpenProblem} (mathematical deficit), and \texttt{MissingFoundation} (ontological deficit).

\subsection{II. The ZFCₜ Architecture}

The library introduces \textbf{ZFCₜ} (ZFC with Temporal depth and Winding), extending the standard Zermelo-Fraenkel axioms into a dynamic, sequential regime.

\subsubsection{Primitive Grammars}

The mathematical systems are encoded as 12-primitive tuples (Synthons):
$\mathcal{S} = \langle D, T, R, P, F, K, G, \Gamma, \Phi, H, S, \Omega \rangle$

\begin{itemize}
    \item \textbf{Dimensionality ($D$):} Transitions from local Euclidean spaces ($D_{\infty}$) to holographic monads ($D_{\odot}$).
    \item \textbf{Polarity ($P$):} Crucial for the "Frobenius Special" condition ($P_{pm\_sym}$), which enables $O_{\infty}$ tier closure.
    \item \textbf{Winding ($\Omega$):} Encodes topological invariants (winding numbers in $\mathbb{Z}$ or $\mathbb{Z}_2$) essential for RH and BSD.
\end{itemize}

\subsubsection{Synthon Encodings}

The file \texttt{Synthon.lean} verifies structural identities across disparate physical and mathematical regimes:

\begin{itemize}
    \item \textbf{P-70 Symmetry:} Higgs, Axion, and Inflaton fields are proved structurally identical (\texttt{higgs = axion = inflaton}) under the $K_{slow}$ (slow-roll) regime.
    \item \textbf{Quantum Gravity ($O_{\infty}$):} Formally identified as holographic Frobenius systems ($D_{\odot} + T_{\odot}$), distinct from the 4D local Yang-Mills quantum target ($D_{\infty}$).
\end{itemize}

\subsection{III. Millennium Barrier Taxonomy}

The core of the scholarship lies in \texttt{Barriers.lean}, where the seven problems are mapped to their structural bottlenecks.

\begin{table}[htbp]
\centering
\small
\rowcolors{1}{white}{white}
\begin{tabularx}{\textwidth}{>{\centering\arraybackslash}X >{\centering\arraybackslash}X >{\centering\arraybackslash}X >{\centering\arraybackslash}X}
\toprule
\rowcolor{gray!25}\textbf{Problem} & \textbf{Barrier Type} & \textbf{Structural signature} & \textbf{Missing Object/Certificate} \\
\midrule
\rowcolors{1}{white}{gray!10}
\textbf{RH} & OpenProblem & $\Phi_c^\mathbb{C}$ & \texttt{ZeroFreeStrip 0} \\
\textbf{YM} & MissingFoundation & $G_{\aleph}, H, \Phi_c$ & \texttt{PathIntegralMeasure G} \\
\textbf{NS} & OpenProblem & $s = 1/2$ gap & \texttt{GlobalRegularityCert} \\
\textbf{BSD} & Mixed (Parallel) & $D_{\odot}, \Omega_Z$ & \texttt{BSDRankCertificate} \\
\textbf{Hodge} & OpenProblem & $D_{\odot}, T_{\odot}$ & \texttt{AlgebraicCycleRep} \\
\textbf{P vs NP} & OpenProblem & $P_{asym} \to P_{pm\_sym}$ & \texttt{CircuitLowerBound} \\
\textbf{OPN} & Mixed (Stacked) & $K_{trap}, \Phi_c$ & \texttt{euler\_opn\_structure} \\
\bottomrule
\end{tabularx}
\end{table}

\subsubsection{Key Formal Theorems:}

\begin{enumerate}
    \item \textbf{\texttt{ym\_is\_unique\_missing\_foundation}}: Proves via exhaustive case analysis that only Yang-Mills requires the construction of an object (\texttt{PathIntegralMeasure}) whose type is currently uninhabitable.
    \item \textbf{\texttt{critical\_scaling\_gap}}: A \texttt{norm\_num} verification in \texttt{NS.lean} confirming the 3D Navier-Stokes critical Sobolev exponent $s=1/2$ sits precisely between the subcritical energy ($s=0$) and supercritical enstrophy ($s=1$).
    \item \textbf{\texttt{rh\_leyang\_correspondence}}: Establishes a structural bridge between the Riemann Hypothesis and the Lee-Yang theorem through the shared \texttt{Phi\_c\_complex} criticality.
\end{enumerate}

\subsection{IV. The Frobenius Tier and Ouroboricity}

The library verifies the \textbf{Crystal address} arithmetic for complex-time path integrals ($6,678,416$) and identifies the \textbf{Frobenius Cliff}. It is machine-checked that $O_{\infty}$ closure (Ouroboricity) is unreachable via tensor products of sub-Frobenius systems --- the $P_{pm\_sym}$ condition is non-synthesizable.

\subsection{V. Conclusion}

\texttt{MILLANKH} provides a meta-mathematical map of the frontier. By formalizing the \emph{reason} for the existence of \texttt{sorry} markers, it transforms human mathematical intuition into machine-verifiable structural constraints, paving the way for targeted foundational expansions in ZFCₜ.


\end{document}
